% =========================================================
% Proof: Cauchy Sequences Are Closed Under Addition and Multiplication
% =========================================================

\newpage
\phantomsection
\label{prf:cauchy-closed-operations-q}

\begin{remark*}[Return]
\hyperref[thm:cauchy-closed-operations-q]{Return to Theorem~\ref*{thm:cauchy-closed-operations-q}.}
\end{remark*}

\begin{theorem*}[Cauchy sequences are closed under addition and multiplication]
If $(x_n)$ and $(y_n)$ are Cauchy sequences of rational numbers, then so
are $(x_n + y_n)$ and $(x_n y_n)$.
\end{theorem*}

\begin{proof}[Proof (Professional Standard)]
\textbf{Addition.}
Let $\varepsilon > 0$.  Since $(x_n)$ and $(y_n)$ are Cauchy, there exist
$N_1, N_2 \in \mathbb{N}$ such that
$|x_m - x_n| < \varepsilon/2$ for all $m,n \geq N_1$ and
$|y_m - y_n| < \varepsilon/2$ for all $m,n \geq N_2$.
Set $N = \max\{N_1, N_2\}$.  For $m, n \geq N$:
\[
|(x_m + y_m) - (x_n + y_n)|
\leq |x_m - x_n| + |y_m - y_n|
< \frac{\varepsilon}{2} + \frac{\varepsilon}{2} = \varepsilon.
\]

\textbf{Multiplication.}
By \hyperref[thm:cauchy-bounded-q]{Theorem~\ref*{thm:cauchy-bounded-q}},
both sequences are bounded; let $\Lambda > 0$ satisfy $|x_n| \leq \Lambda$
and $|y_n| \leq \Lambda$ for all $n$.
Let $\varepsilon > 0$.  Choose $N_1, N_2$ such that
$|x_m - x_n| < \varepsilon/(2\Lambda)$ for $m,n \geq N_1$ and
$|y_m - y_n| < \varepsilon/(2\Lambda)$ for $m,n \geq N_2$.
Set $N = \max\{N_1, N_2\}$.  For $m, n \geq N$:
\[
|x_m y_m - x_n y_n|
= |x_m y_m - x_n y_m + x_n y_m - x_n y_n|
\leq |y_m||x_m - x_n| + |x_n||y_m - y_n|
< \Lambda \cdot \frac{\varepsilon}{2\Lambda} + \Lambda \cdot \frac{\varepsilon}{2\Lambda}
= \varepsilon.
\]
\end{proof}

\begin{proof}[Proof (Detailed Learning)]
\textbf{Part I: Addition.}

\textbf{Step 1: Set up the $\varepsilon/2$ split.}
Let $\varepsilon > 0$.  We need $|(x_m+y_m)-(x_n+y_n)| < \varepsilon$
for all sufficiently large $m,n$.  We will achieve this by keeping both
$|x_m-x_n|$ and $|y_m-y_n|$ smaller than $\varepsilon/2$.

\textbf{Step 2: Extract witnesses.}
Since $(x_n)$ is Cauchy, $\exists N_1$ such that $|x_m-x_n|<\varepsilon/2$
for all $m,n \geq N_1$.
Since $(y_n)$ is Cauchy, $\exists N_2$ such that $|y_m-y_n|<\varepsilon/2$
for all $m,n \geq N_2$.

\textbf{Step 3: Combine and apply the triangle inequality.}
Set $N = \max\{N_1, N_2\}$.  For all $m, n \geq N$:
\[
|(x_m + y_m) - (x_n + y_n)|
= |(x_m - x_n) + (y_m - y_n)|
\leq |x_m - x_n| + |y_m - y_n|
< \frac{\varepsilon}{2} + \frac{\varepsilon}{2} = \varepsilon.
\]

\textbf{Part II: Multiplication.}

\textbf{Step 4: Invoke boundedness.}
By \hyperref[thm:cauchy-bounded-q]{Theorem~\ref*{thm:cauchy-bounded-q}},
$(x_n)$ and $(y_n)$ are bounded.  Let $\Lambda_1, \Lambda_2 > 0$ satisfy
$|x_n| \leq \Lambda_1$ and $|y_n| \leq \Lambda_2$ for all $n$.
Set $\Lambda = \max\{\Lambda_1, \Lambda_2\} > 0$, so both $|x_n| \leq \Lambda$
and $|y_n| \leq \Lambda$ for all $n$.

\textbf{Step 5: Algebraically separate the two sources of error.}
Add and subtract the cross-term $x_n y_m$:
\[
x_m y_m - x_n y_n
= (x_m - x_n)y_m + x_n(y_m - y_n).
\]
By the triangle inequality:
\[
|x_m y_m - x_n y_n|
\leq |y_m||x_m - x_n| + |x_n||y_m - y_n|
\leq \Lambda|x_m - x_n| + \Lambda|y_m - y_n|.
\]

\textbf{Step 6: Choose the Cauchy witnesses to compensate for $\Lambda$.}
Let $\varepsilon > 0$.  Choose $N_1$ such that $|x_m-x_n| < \varepsilon/(2\Lambda)$
for all $m, n \geq N_1$, and $N_2$ such that $|y_m-y_n| < \varepsilon/(2\Lambda)$
for all $m, n \geq N_2$.

\textbf{Step 7: Combine.}
Set $N = \max\{N_1, N_2\}$.  For all $m, n \geq N$:
\[
|x_m y_m - x_n y_n|
\leq \Lambda \cdot \frac{\varepsilon}{2\Lambda} + \Lambda \cdot \frac{\varepsilon}{2\Lambda}
= \frac{\varepsilon}{2} + \frac{\varepsilon}{2} = \varepsilon.
\]

\textbf{Step 8: Conclude.}
Both $(x_n + y_n)$ and $(x_n y_n)$ satisfy the Cauchy condition.
\end{proof}

\begin{remark*}[Proof structure]
Both parts are $\varepsilon$-splitting arguments.  The addition case is a
direct two-term triangle inequality.  The multiplication case adds one
layer of difficulty: the product does not split cleanly, so one must first
algebraically separate the two sources of error (the $x$-difference and
the $y$-difference) using the cross-term trick, then compensate for the
bounds $\Lambda$ by scaling the tolerance down to $\varepsilon/(2\Lambda)$
before extracting the Cauchy witnesses.  This $\Lambda$-compensation pattern
recurs throughout analysis wherever products of sequences appear.
\end{remark*}

\begin{remark*}[Dependencies]
\hyperref[def:cauchy-rational]{Cauchy sequence in $\mathbb{Q}$},
\hyperref[thm:cauchy-bounded-q]{Cauchy sequences are bounded},
the triangle inequality in $\mathbb{Q}$.
\end{remark*}
